\documentclass[11pt, a4paper]{article}

% --- Paquets ---
\usepackage[utf8]{inputenc}
\usepackage[T1]{fontenc}
\usepackage[french]{babel}
\usepackage[margin=2.5cm]{geometry}
\usepackage{hyperref}
\usepackage{enumitem}
\usepackage{booktabs}
\usepackage{listings}
\usepackage{xcolor}

% --- Configuration Code ---
\lstset{
    language=Python,
    basicstyle=\ttfamily\small,
    keywordstyle=\color{blue},
    stringstyle=\color{red},
    commentstyle=\color{gray},
    breaklines=true,
    frame=single,
}

\title{\textbf{Rapport de recherche : Bibliothèque ByzFL} \\ \large Reproductibilité, Données Non-IID et Limites Théoriques}
\author{ABBES Samy \& AREZKI Célia}
\date{Semaine du 11 Mai 2026}

\begin{document}

\maketitle

\section{Introduction}
Suite à la réunion du 11/05/2026, cette semaine a été consacrée à la professionnalisation de notre environnement de travail et à l'exploration de scénarios plus complexes avec les jeux de données MNIST et FashionMNIST. Nous avons étudié l'impact d'une distribution hétérogène (Non-IID) sur les performances des agrégateurs et avons testé empiriquement les limites mathématiques de tolérance aux clients byzantins.

\section{Réorganisation du dépôt et Reproductibilité}
Afin de garantir un environnement de recherche robuste et collaboratif, nous avons restructuré notre dépôt Git :
\begin{itemize}
    \item \textbf{Séparation des dossiers} : Les codes sources (\texttt{code/}), les données brutes (\texttt{data/}), et les rapports (\texttt{rapports/}) sont désormais isolés. Cela évite les téléchargements redondants de jeux de données et clarifie l'arborescence.
    \item \textbf{Environnement Virtuel} : Nous avons configuré l'environnement avec le gestionnaire \texttt{uv}.
    \item \textbf{Fichiers de configuration} : La mise en place d'un \texttt{README.md}, d'un fichier \texttt{requirements.txt} et d'un \texttt{uv.lock} permet désormais à quiconque de reproduire nos expériences à l'identique, depuis n'importe quelle machine, en figeant les versions exactes des dépendances.
\end{itemize}

\section{Impact de la distribution des données (FashionMNIST)}
Nous avons remplacé MNIST par FashionMNIST, un jeu de données plus complexe, et avons comparé l'apprentissage sur 100 cycles (rounds) entre une distribution homogène (IID) et une distribution hétérogène (Non-IID via une loi de Dirichlet).

\begin{table}[h]
\centering
\begin{tabular}{@{}llcc@{}}
\toprule
\textbf{Scénario} & \textbf{Agrégateur (Attaque)} & \textbf{Précision (IID)} & \textbf{Précision (Non-IID)} \\ \midrule
clean\_average     & Average (Aucune)              & 0.7065 & 0.7210 \\
signflip\_average  & Average (SignFlipping)        & 0.6445 & 0.6655 \\
signflip\_median   & Median (SignFlipping)         & \textbf{0.6950} & \textbf{0.4585} \\
ipm\_trmean        & TrMean (IPM)                  & \textbf{0.7158} & \textbf{0.5071} \\
alie\_krum         & Krum (ALIE)                   & \textbf{0.7067} & \textbf{0.6157} \\ \bottomrule
\end{tabular}
\caption{Comparaison des précisions finales sur FashionMNIST (100 rounds)}
\end{table}

\subsection{Analyse des résultats}
L'expérience met en évidence une vulnérabilité majeure des agrégateurs robustes face aux données Non-IID. En IID, \textit{Median}, \textit{TrMean} et \textit{Krum} résistent parfaitement aux attaques et maintiennent une précision d'environ 70\%. En revanche, avec des données Non-IID, les performances s'effondrent (chute à 45\% pour la médiane et 50\% pour TrMean). Cela s'explique par le fait que les agrégateurs peinent à distinguer une variation de gradient due à une attaque d'une variation due à une différence légitime de données locales.


\section{Preuve empirique du théorème de robustesse ($2f \ge n$)}
Nous avons aussi conçu une expérience visant à briser mathématiquement l'agrégateur géométrique Krum en utilisant l'attaque \textit{SignFlipping} et en modifiant progressivement le ratio de clients byzantins ($f$) par rapport au nombre total de clients ($n$).

\textbf{Résultats obtenus :}
\begin{itemize}
    \item \textbf{5 Honnêtes vs 1 Byzantin ($f=1, n=6$)} : Précision finale de 77.64\%. Krum filtre facilement le gradient malveillant.
    \item \textbf{4 Honnêtes vs 2 Byzantins ($f=2, n=6$)} : Précision finale de 74.78\%. La majorité honnête est maintenue, le modèle converge.
    \item \textbf{3 Honnêtes vs 3 Byzantins ($f=3, n=6$)} : \textbf{Erreur d'exécution}.
\end{itemize}

\subsection{Explication de l'erreur ByzFL}
Lors du dernier test, le programme s'est interrompu avec l'erreur : \texttt{ValueError: Cannot tolerate 2f >= n. Got f=3, n=6}. 
Les algorithmes d'agrégation robuste reposent sur l'hypothèse fondamentale qu'il existe une majorité absolue de clients honnêtes. Dès lors que $2f \ge n$ (ici 50\% de byzantins), les attaquants ont la capacité de dicter la direction du modèle global. Le framework détecte cette faille structurelle et préfère lever une exception de sécurité plutôt que d'empoisonner le modèle.


\end{document}
